\section{Open Models Go Vertical}
\sectionkicker{OPEN VERTICAL}
\begin{wideflow}
{\Large\bfseries オープンモデルの3方向}\par
\smallskip
{\color{SurveyMuted}9月7日から10日の3件の公開を、オンデバイス用途、翻訳特化、画像・動画対応の3方向から、日付とライセンスの精度とともに示す。}\par
\end{wideflow}

OpenBMBのMiniCPM5-2Bは、オンデバイスやローカル実行を狙う20億級の密な言語モデルである\cite{minicpm}。Llama系の標準的な構成で42層、コンテキスト長は約13万トークン、Apache-2.0ライセンスで公開される。同規模比較での平均53.9、LiveCodeBench 69.1、AIME 86.5、MATH-500 94.6、SWE-bench Verified 46.4、GAIA 88.7、tau2 Telecom 97.1、BFCL 66.6といった値はモデルカードの数値であり、外部の再現はない。vLLM・SGLang・llama.cpp・Transformersでフォークなしに動作し、ツール呼び出しの受け口やUltraDataファミリーの公開、強化学習とモデル統合の改善が添えられるが、量子化やチップごとの性能はベンダー側の説明にとどめる。9月7日という日付はレシピ公開や周辺記録からの帰属であり、時刻レベルの断定はしない。

9月10日、CohereはNorth Small Translate 1.0を公開した\cite{north}。翻訳に特化したMoEモデルで全体218B、動作パラメーター25B、入出力とも16Kトークン、50以上の言語をサポートするとされる。ライセンスは研究用途のCC BY-NC 4.0で、商用利用には別途Vaultライセンスが必要とされる。RWSとの協業、Hugging Faceリポジトリやドキュメント、テクニカルレポートが添えられる。WMT26全言語で83.60（エージェント型で84.36）、DeepL・Google翻訳・Gemma・GLM・Mistralとの対比が示されるが、評価にはGPT-5.6 Solが評価用モデルとして使われ、サンプリングや言語別の詳細表は今回たどっていない。スループットやコストの内訳（Gemma 31B比1.4倍、長文脈48.9、タスクあたりコストなど）は社内計測として受け止める。

InclusionAIのLing-3.0-flash-VLは、Ling-3.0-flashをベースに画像・動画理解を追加したネイティブマルチモーダルモデルである\cite{ling}。全体124B、トークンあたり5.5Bが動作し、最大256Kトークンのコンテキスト、画像・動画入力に対応し、MITライセンスで公開されるとされる。Artificial Analysis Intelligence Indexで42（ベースのflashから4ポイント増）、理解・推論・行動の3軸評価やSGLang・vLLM向け手順が添えられるが、数値はモデルカードの報告であり独立した再現はない。リポジトリ作成は9月4日、FP8版は9月8日、モデル紹介ページの掲載は9月10日という周辺記録で押さえ、時刻レベルの断定はしない。動画対応の有効性はベンダー側の説明として扱う。

\begin{claimboundary}[オープンモデルの境界]
数値はモデルカードの報告であり、外部の再現はない。日付は帰属の精度ごと扱い、時刻レベルの断定はしない。ライセンス範囲を事実と混同しない。
\end{claimboundary}

\begin{communitynote}[X / Community observation --- context only]
オンデバイスでの試行や速度の報告が、9月7日の公式発信の後に独立した手で重なる様子が観察される\cite{w37x-minicpm-official,w37x-minicpm-ind}。うかがえるのはローカル実行への関心であり、性能や対応チップの根拠にはならない。
\end{communitynote}
